\documentclass[./测控三级项目.tex]{subfiles} % 使用 subfiles 文档类，指定主文档

\begin{document}
	\subsection{题目重述}
		\begin{compactenum}
			\item 设计抑制题6[1]e)中的中频信号的滤波器电路，用Multisim对滤波器的频率特性进行仿真。
			\item 搭建实测电路，分别为上述低通实测电路引入一个 20mV 正弦信号，调整信号频率，在 0.1Hz-1MHz 范围内测试 30 个以上关键频率点的输出波形，记录输出信号幅值变化，用 MATLAB 描点绘制 Bode 图，用术语描述滤波器的频率特性。
		\end{compactenum}
	\subsection{题目分析与实施方案}
		\subsubsection{第一问}
			第一问要求抑制题6[1]e)中的中频信号，选用了Sally-Key结构二阶有源带阻滤波器。Sally-Key结构二阶有源带阻滤波器是由一个双T网络和一个经由顶部电容提供正反馈的增益单元组成。因此考虑在本题使用这个电路。 \par
			电路确定后，通过计算有源带阻网络的传递函数，得到对应相频特性函数。将需要抑制的频率代入特性曲线的截止频率，便可以推算出对应电路参数，并根据实际情况，选用相关精度的标准阻值电阻。 \par
			然后，使用 Multisim 软件仿真对应电路，将题6（1）[e] 中的输出信号输入搭建的Sally-Key结构二阶有源带阻滤波器上，并连接示波器。
		\subsubsection{第二问}
			根据第一问搭建的电路，向Sallen-Key结构二阶有源带阻滤波器等对数频率，分别测试输出信号的幅值变化三十组。\par 
		使用 MATLAB 的 plot 函数，描点绘制出测试得到的参数，以得到相频伯德图曲线。\par 
		根据绘制的伯德图曲线，分析该Sallen-Key结构二阶有源带阻滤波器的频率特性。
	\subsection{理论分析与仿真}
		\subsubsection{第一问}
			图\ref{fig:Sallen-Key结构二阶有源带阻滤波器电路原理图}是一个Sallen-Key结构二阶有源带阻滤波器电路。
			\begin{figure}[h]
				\centering
				\includegraphics[width=0.45\textwidth]{Sallen-Key结构二阶有源带阻滤波器电路原理图}
				\caption{Sallen-Key结构二阶有源带阻滤波器电路原理图}
				\label{fig:Sallen-Key结构二阶有源带阻滤波器电路原理图}
			\end{figure}\par
			根据电路计算可得Sallen-Key结构二阶有源带通滤波器电路的传递函数。令$s=j\omega$，并将其表示为归一化形式，可得其归一化频率$\omega_0$、归一化增益$H_0$和归一化品质因数$Q$，为了方便参数计算，以下采用等值元件法简化。令$C_1=C_2=C$，则归一化方程组可化为：
			\begin{gather}
				\begin{dcases}
					H_0=1+\frac{R_A}{R_B}\\
					\omega_0 =\frac{1}{RC}		\\
					Q =\frac{R_A}{2R_A-2R_B}
					\label{eq:Sallen-Key结构二阶有源带阻滤波器截止频率}
				\end{dcases}
			\end{gather}\par
			以上理论完备，在Multisim软件中进行仿真实验。\par
			由于需要滤除的中频信号为500Hz，需要保留的低频信号为200Hz，高频信号为1kHz。因此，设置中心频率为$f_0=500$Hz，归一化品质因数$Q=0.35$。\par
			设置电容容值为$C=100$nF(Standeard Value,$5\%$)\par
			根据式\ref{eq:Sallen-Key结构二阶有源带阻滤波器截止频率}，电阻阻值$R=3473.04559\Omega$，与标准电阻值表比对后，实际选用电阻阻值为$R=3.48$k$\Omega$(Standeard Value,$1\%$)的电阻。\par
			将要求的信号接入该电路，得到仿真电路图，如图\ref{fig:Sallen-Key结构二阶有源带阻滤波器仿真原理图}所示。
			\begin{figure}[H]
				\centering
				\includegraphics[width=0.45\textwidth]{./14.第十四题示波器现象/14仿真原理图}
				\caption{Sallen-Key结构二阶有源带阻滤波器仿真原理图}
				\label{fig:Sallen-Key结构二阶有源带阻滤波器仿真原理图}
			\end{figure}\par
			对输入和输出信号做傅里叶变换，得到对应频谱曲线，如图\ref{fig:Sallen-Key结构二阶有源带阻滤波器频谱曲线}所示。
			\begin{figure}[h]
				\centering
				\includegraphics[width=0.45\textwidth]{./14.第十四题示波器现象/14仿真频率曲线}
				\caption{Sallen-Key结构二阶有源带阻滤波器频谱曲线}
				\label{fig:Sallen-Key结构二阶有源带阻滤波器频谱曲线}
			\end{figure}\par
			这说明图\ref{fig:Sallen-Key结构二阶有源带阻滤波器电路原理图}的带阻滤波器可以对中频信号进行滤波操作。
	\subsection{实际电路测试结果与数据分析}
	\subsubsection{第二问}
	按题目要求引入幅值为20mV的正弦信号，并调整信号频率，在0.1Hz-1MHz区间测试30个以上关键频率点，测试结果如表\ref{tab:不同频率下Sallen-Key结构二阶有源带阻滤波器增益表}所示。
	\begin{table}[h!]
		\centering
		\caption{不同频率下Sallen-Key结构二阶有源带阻滤波器增益表}
		\label{tab:不同频率下Sallen-Key结构二阶有源带阻滤波器增益表}
		\begin{tabular}{|c|c|c|c|c|c|c|c|}
			\hline
			频率（Hz） & 增益  & 频率（Hz） & 增益 & 频率（Hz） & 增益 & 频率（Hz） & 增益\\
			\hline
			0.1&1.57&10&1.57&1000&1.4&100000&1.57\\\hline
			0.3&1.57&30&1.57&3000&1.55&300000&1.57\\\hline
			0.5&1.57&50&1.56&5000&1.56&500000&1.57\\\hline
			0.7&1.57&70&1.55&7000&1.57&700000&1.57\\\hline
			0.9&1.57&90&1.54&9000&1.57&900000&1.57\\\hline
			1&1.57&100&1.54&10000&1.57&1000000&1.57\\\hline
			3&1.57&300&1.11&30000&1.57&&\\\hline
			5&1.57&500&0.32&50000&1.57&&\\\hline
			7&1.57&700&1.11&70000&1.57&&\\\hline
			9&1.57&900&1.35&90000&1.57&&\\\hline
		\end{tabular}
	\end{table}	\par
	将表\ref{tab:不同频率下Sallen-Key结构二阶有源带阻滤波器增益表}中的数据输入MATLAB中进行Bode图生成，其图像如下图\ref{fig:Sallen-Key结构二阶有源带阻滤波器测试伯德图}所示。
	\begin{figure}[h]
		\centering
		\includegraphics[width=0.7\textwidth]{./14.第十四题示波器现象/Sallen-Key结构二阶有源带阻滤波器测试伯德图}
		\caption{测试生成Sallen-Key结构二阶有源带阻滤波器伯德图}
		\label{fig:Sallen-Key结构二阶有源带阻滤波器测试伯德图}
	\end{figure}\par
	根据实际选取电阻电容值参数，可以直接得到对应传递函数。利用MATLAB软件，得到该传递函数的伯德图如图\ref{fig:Sallen-Key结构二阶有源带阻滤波器实际伯德图}所示。
	\begin{figure}[H]
		\centering
		\includegraphics[width=0.7\textwidth]{./14.第十四题示波器现象/Sallen-Key结构二阶有源带阻滤波器实际伯德图}
		\caption{理论生成Sallen-Key结构二阶有源带阻滤波器伯德图}
		\label{fig:Sallen-Key结构二阶有源带阻滤波器实际伯德图}
	\end{figure}\par
	对比图\ref{fig:Sallen-Key结构二阶有源带阻滤波器测试伯德图}和图\ref{fig:Sallen-Key结构二阶有源带阻滤波器实际伯德图}，中的相频曲线，可以发现除了曲线前者不如后者光滑外，其余基本一致，这说明实际测试结果符合理论推断。\par
	同时也注意到，在100Hz-2kHz范围内，增益大幅变化，从1下降到0，再上升回到了1的位置，这说明带阻滤波器在该频段才具有过滤信号的特点。
	\subsection{结论与收获}
	\textbf{游寅龙}:\par
	理论分析层面:查找了二阶有源带阻滤波器的电路原理图，深入理解了Sallen - Key结构二阶有源带阻滤波器的电路原理。\par
	动手实践方面:能熟练使用信号发生器产生20mV正弦信号，并能运用示波器对结果进行观测和记录，提高了动手能力。\par
	精神层面: 发现了技术工具的便捷性，开阔了眼界，减少了不必要的时间浪费。\par
	\textbf{郭嘉媛}:\par
	理论分析层面:了解其基本的电路组成和工作原理，能通过推导和分析传递函数，理解其在不同频率下的增益、相移等特性变化规律，增强了总结能力。\par
	动手实践层面:完成了Multsimn 仿真的搭建，能够根据理论分析实时调整搭建的电路，增强了独立解决问题的能力并丰富了解决问题的手段。\par
	\textbf{郑嘉宁}:\par
	理论分析层面:理解了滤波器截止频率在电路设计中的关键作用，能够根据不同的题目要求计算相应的截止频率，知识掌握更加熟练。\par
	动手实践层面:完成了对实测电路的搭建，能够根据要求调试电路，焊接能力有很大提升。\par
	精神层面:学到了有问题提出来，大家一起解决，而不是一个人埋头苦干，但没有效率。从而加快了工作的效率。
	
\end{document}